\documentclass[10pt,a4paper]{article}
\usepackage[margin=1.55cm]{geometry}
\usepackage{xcolor}
\usepackage{amsmath,amssymb}
\usepackage{graphicx}
\usepackage{enumitem}
\usepackage{fontspec}
\usepackage[CJKmath=true]{xeCJK}
% CJK main font is resolved at COMPILE time on the actual machine, not baked in at
% generation time, so a .tex generated on one OS still renders on another (the older
% Python-side probe hardcoded one OS's font name into the file and broke when moved).
% Walk a cross-platform priority list and use the first font that exists; AutoFakeBold
% + AutoFakeSlant give bold/italic CJK even on faces that ship no bold/italic cut.
\newcommand{\setCJKto}[1]{\setCJKmainfont{#1}[AutoFakeBold=true,AutoFakeSlant=0.2]}
\IfFontExistsTF{Noto Sans CJK SC}{\setCJKto{Noto Sans CJK SC}}{%        % Linux (fonts-noto-cjk); cross-platform if installed
 \IfFontExistsTF{Source Han Sans SC}{\setCJKto{Source Han Sans SC}}{%   % Adobe 思源黑体; cross-platform
  \IfFontExistsTF{PingFang SC}{\setCJKto{PingFang SC}}{%               % macOS 10.11+
   \IfFontExistsTF{Hiragino Sans GB}{\setCJKto{Hiragino Sans GB}}{%    % older macOS
    \IfFontExistsTF{Microsoft YaHei}{\setCJKto{Microsoft YaHei}}{%     % Windows
     \IfFontExistsTF{WenQuanYi Micro Hei}{\setCJKto{WenQuanYi Micro Hei}}{% % older Linux
      \setCJKto{Songti SC}}}}}}}                                       % macOS bottom fallback
\definecolor{cblue}{HTML}{4C72B0}
\setlength{\parskip}{3pt}\setlength{\parindent}{0pt}
\setlist[enumerate]{topsep=2pt,itemsep=2pt,parsep=1pt,partopsep=0pt}
\setlist[itemize]{topsep=2pt,itemsep=2pt,parsep=1pt,partopsep=0pt}
% Let prose-embedded inline math break across lines and absorb small overflows, so simple
% formulas inserted mid-sentence stop running past the right margin.
\tolerance=1200
\emergencystretch=2.5em
\relpenalty=20
\binoppenalty=40
\hfuzz=1pt
\begin{document}

\section*{用于内生高频 VO2 开关前沿的 Kinetics-Clock PINNs}
\textbf{方法名称：} Kinetics-Clock Physics-Informed Neural Network (KC-PINN)
\section*{研究动机}
结构性问题在于，内生开关规律会把大部分相变压缩到物理时间中一个很短、而且会在空间中移动的区段。在二维 VO2 器件中，不同位置会以不同速率经历相变，而电压、温度和相态仍然紧密耦合。只使用全局物理时间的网络必须在同一坐标中同时表示尖锐前沿和平滑背景，因此会暴露 arxiv:2602.19265v1 所强调的高频和残差条件问题，以及 arxiv:2605.04502v1 的自适应频谱诊断所指出的问题。

最接近的架构先例 NeuSA（arxiv:2509.04966v2）为典型波动问题加入了全局频谱状态和因果演化，但没有从材料动力学构造空间局部时钟。openalex:W4403555504 中的能量耗散相场物理信息神经网络（PINN）通过求解静态稳态问题来避开传播失败，因此不能表示由脉冲产生的事件时间或移动拓扑。Kinetics-Clock PINN（KC-PINN）则使用有限解析 Gaussian 基及其误差函数原语构造局部时钟，不进行数值时间求积。原始电-热-相态方程仍然是权威，所有一阶和二阶导数都会通过学习到的映射变换回去，其中包括跨器件变化的时钟所产生的全部空间 Hessian 项和交叉项。

三项近期工作使这个狭窄缺口现在可以处理。arxiv:2605.04502v1 分离了自适应频谱 PINN 中的梯度缩放效应，arxiv:2602.19265v1 汇集了高频失败机制，semanticscholar:a1448fd2600739b73894f3b8d467bf9a9e9630b2 则揭示了亚纳秒 VO2 电-热-相态拓扑这一相关物理状态。指定的 FEniCSx 模型可以提供受控合成轨迹，自动微分可以计算完整的局部坐标变换，而不是近似或遗漏其中的交叉项。因此机会不只是增加频谱容量，而是在检索边界内的新颖性主张下，检验与内生动力学对齐的坐标能否移除这个特定的表示瓶颈。

NeuSA 使用由神经常微分方程演化的截断全局频谱状态和因果物理约束来表示典型初值问题。它没有学习与材料相变速率规律相连的空间局部单调时钟，也没有通过该时钟变换耦合电-热-相态器件偏微分方程（PDE）。它的全局频谱演化坐标缺少定义局部材料时间所需的内生、空间异质动力学信号。自适应频谱研究改变 Fourier 频率和初始条件门，以诊断刚性非线性常微分方程中的梯度缩放，但没有变换二维耦合器件 PDE 的时间，也没有保留空间变化映射所诱导的链式法则项。其诊断设置中没有空间分辨的相态动力学右端项，因而无法据此构造并检验内生局部时钟。能量嵌入式相场 PINN 通过在静态形式中强制能量耗散来预测稳态铁电微结构。它用稳态变分目标替代轨迹问题来处理动态 PINN 的传播失败，而没有保留脉冲驱动的动态演化、开关时间和移动相前沿拓扑。原位 VO2 研究分辨了电控导电拓扑和亚纳秒电-热-相态动力学。它的贡献是物理测量和机制辨别，而不是可微 PDE 求解器或从测得动力学导出的训练坐标，因此没有提供用于求解器件控制场的 PINN 坐标机制，也没有在合成数值证据上分离动力学特异的表示价值。

弥合这一缺口后，一个二维 PINN 可以共同表示由脉冲产生的事件时间、移动前沿拓扑和电-热能量闭合，而不把底层方程替换为静态代理。这会把 semanticscholar:a1448fd2600739b73894f3b8d467bf9a9e9630b2 所揭示的动态目标状态连接到可在完整留出器件或协议上评估的求解器，并直接检验当动力学在空间上不均匀时，arxiv:2509.04966v2 的全局频谱补救是否不足。它还会形成一个可复用的条件性表述：当局部相变速率幅值具有信息性且变换完整时，坐标自适应可归因于材料动力学，而不是一般的额外容量。
\section*{方法}
\subsection*{M0\_background}
\emph{冻结所有分支共用的二维电-热-相态方程、脉冲合同、条件和残差尺度。}
\begin{enumerate}[leftmargin=*,start=1]
  \item 把现有二维器件模型整理成只读的 `PhysicalContract`。其中保存网格与标签、材料区域、脉冲分段及每个跳变处的左右单侧值、带明确输入的材料函数、初始状态、边界与界面条件、单位，以及每种无量纲换算的正向和逆向规则；后续步骤必须读取同一版本。 【作者需决定：指定权威 FEniCSx 模型快照，并批准完整物理合同清单、参考量和残差尺度规则。】 由该合同推导无量纲电、热和相态算子，不增加控制项。只用训练案例为每类内部、初始、边界、本构和脉冲界面残差建立有名称的尺度，并在验证评分前冻结。输出物理算子、脉冲界面表、单位换算和残差尺度规范。
\par\emph{为什么：} 这使局部时钟成为表示变换，而不是对物理模型的替换，并固定用于评定耦合残差和能量闭合的量。
\end{enumerate}
\subsection*{M1\_analytic\_clock}
\emph{由空间系数网络和解析 Gaussian 误差函数原语构造严格单调的局部时钟。}
\begin{enumerate}[leftmargin=*,start=2]
  \item 建立 `ClockBasisSpec`，记录无量纲时间区间、固定高斯中心 \(`c_{j}`\)、公共间距 `w`、浮点类型、`eps\_dtype`、严格为正的 `kref`、由此得到的正下限 `kfloor`，以及候选集合 `\{8,16,32\}`。 【作者需决定：在推导 \(`kfloor=sqrt(eps_{dtype})\times kref`\) 前，锁定 \(`c_{j}`\)、`w`、浮点类型，以及严格为正的 `kref` 的单位、来源和数值。】 使用只读取空间位置的网络 \(`a_{psi}(x)`\)，使其恰好输出 \(`M_{tau}`\) 个系数，且不得读取脉冲或任何预测场。 【作者需决定：声明系数网络的输入归一化、层结构、激活函数，以及选择 \(`M_{tau}`\) 的单一训练---验证评分；默认 \(`M_{tau}=16`\)，完全并列时选较小候选。】 用 softplus 使每个系数为正，求值固定高斯函数及其解析误差函数累积形式，并严格按 E1 形成 `tau\_psi`。时间与空间导数由解析基函数和对空间网络的自动微分得到，绝不使用有限时间积分。用 `ClockState` 返回系数、基函数值、时钟、全部所需时钟导数、选定的 \(`M_{tau}`\) 和冻结的基规范。
\noindent\emph{时钟由高斯基函数的解析误差函数原语直接计算，不使用有限时间求积。固定结点 \(`c_{j}`\)、间距 `w`、预先声明的正参考速率 `kref>0`、`kfloor=sqrt(eps\_dtype)*kref`，并仅在训练---验证集上从 \(`M_{tau}\) in \{8,16,32\}` 选择，可使 `tau\_psi(x,0)=0` 且实现中的 \(`partial_{t}\) tau\_psi>0` 精确成立。}
\par\noindent{\small\ttfamily [eq~1, raw LaTeX, not typeset] \textbackslash\{\}tau\_\textbackslash\{\}psi(\textbackslash\{\}mathbf\{x\},t)=k\_\{\textbackslash\{\}mathrm\{floor\}\}t+\textbackslash\{\}sum\_\{j=1\}\textasciicircum{}\{M\_\textbackslash\{\}tau\}\textbackslash\{\}operatorname\{softplus\}(a\_\{\textbackslash\{\}psi,j\}(\textbackslash\{\}mathbf\{x\}))B\_j(t),\textbackslash\{\}quad b\_j(t)=e\textasciicircum{}\{-((t-c\_j)/w)\textasciicircum{}2\},\textbackslash\{\}quad B\_j(t)=\textbackslash\{\}frac\{\textbackslash\{\}sqrt\{\textbackslash\{\}pi\}w\}\{2\}\textbackslash\{\}left[\textbackslash\{\}operatorname\{erf\}\textbackslash\{\}!\textbackslash\{\}left(\textbackslash\{\}frac\{t-c\_j\}\{w\}\textbackslash\{\}right)+\textbackslash\{\}operatorname\{erf\}\textbackslash\{\}!\textbackslash\{\}left(\textbackslash\{\}frac\{c\_j\}\{w\}\textbackslash\{\}right)\textbackslash\{\}right],\textbackslash\{\}quad k\_\{\textbackslash\{\}mathrm\{ref\}\}>0}
\par\emph{为什么：} 该解析构造使 `tau\_psi(x,0) = 0` 和 `tau\_psi` 的时间导数严格为正这两个条件精确成立，并且不使用有限求积。系数网络不读取任何预测场，因此前向图是无环的。
\end{enumerate}
\subsection*{M2\_fields\_pullback\_residuals}
\emph{预测耦合场，完成考虑脉冲界面的完整导数变换，并计算未改变的物理残差和直接动力学-时钟残差。}
\begin{enumerate}[leftmargin=*,start=3]
  \item 把脉冲协议保存为 `PulseTrace`：每个光滑分段返回脉冲值及其一、二阶时间导数，每个跳变返回独立的左、右单侧值和稳定的 \(`interface_{id}`\)。将位置、`tau\_psi` 和当前单侧脉冲值输入同一个三输出解头 \(`h_{theta}`\)。 【作者需决定：声明容量匹配解头的共享主干、层结构、激活函数、输入归一化，以及电压、温度和相态的输出参数化。】 把输出保存为 \(`V_{hat}`\)、\(`T_{hat}`\) 和 \(`xi_{hat}`\)。由电压的物理空间梯度得到电场，用冻结电导率得到电流密度，再由电流密度与电场的内积得到焦耳热；不得用学习得到的附加头替代这些确定性关系。用 `FieldState` 返回三个场、这些派生量、脉冲记录、时钟引用和物理合同引用。
\par\emph{为什么：} 一个共享时钟保留电学到 Joule 热、再到热学和相态的因果链。已知单侧脉冲值只进入解头，因此时钟计算图仍然无环。
  \item 为每个标量场建立 `DerivativeBundle`。在脉冲光滑分段内部，用自动微分得到相对于位置、`tau\_psi` 和每个脉冲分量的导数，并保留 E3 与 E4 所需的全部混合导数和二阶导数。将它们与脉冲的一、二阶时间导数及精确时钟导数组合，严格形成 E3 的完整一阶和二阶物理时间导数。再严格按 E4 形成完整 \(`2\times 2`\) 物理空间黑塞矩阵，包括两项位置---时钟混合项、时钟梯度外积和时钟黑塞矩阵。黑塞矩阵始终保留为矩阵；拉普拉斯量只能由其迹得到，散度只能对完整待求导量形成。遇到脉冲跳变时分别求值左右迹，绝不跨跳变求导。返回全部导数及对应的脉冲分段和侧别标签。
\noindent\emph{对每个标量场 `f=h\_theta(x,tau\_psi,p(t))`，在每个脉冲光滑开区间内计算完整的一阶和二阶时间拉回；脉冲跳变作为界面处理，绝不对跳变本身求导。}
\par\noindent{\small\ttfamily [eq~2, raw LaTeX, not typeset] \textbackslash\{\}partial\_t f=f\_\textbackslash\{\}tau\textbackslash\{\}tau\_t+f\_p\textbackslash\{\}!\textbackslash\{\}cdot\textbackslash\{\}dot p,\textbackslash\{\}quad \textbackslash\{\}partial\_\{tt\}f=f\_\{\textbackslash\{\}tau\textbackslash\{\}tau\}\textbackslash\{\}tau\_t\textasciicircum{}2+f\_\textbackslash\{\}tau\textbackslash\{\}tau\_\{tt\}+2(f\_\{\textbackslash\{\}tau p\}\textbackslash\{\}!\textbackslash\{\}cdot\textbackslash\{\}dot p)\textbackslash\{\}tau\_t+f\_\{pp\}[\textbackslash\{\}dot p,\textbackslash\{\}dot p]+f\_p\textbackslash\{\}!\textbackslash\{\}cdot\textbackslash\{\}ddot p}
\noindent\emph{先构造完整物理空间 Hessian，包括两项混合缩并、时钟梯度外积和时钟 Hessian，再取迹或散度；Hessian 与 Laplacian 不混用。}
\par\noindent{\small\ttfamily [eq~3, raw LaTeX, not typeset] D\_\{\textbackslash\{\}mathbf\{x\}\}\textasciicircum{}2f=f\_\{\textbackslash\{\}mathbf\{x\}\textbackslash\{\}mathbf\{x\}\}+f\_\{\textbackslash\{\}mathbf\{x\}\textbackslash\{\}tau\}\textbackslash\{\}otimes\textbackslash\{\}nabla\textbackslash\{\}tau\_\textbackslash\{\}psi+\textbackslash\{\}nabla\textbackslash\{\}tau\_\textbackslash\{\}psi\textbackslash\{\}otimes f\_\{\textbackslash\{\}tau\textbackslash\{\}mathbf\{x\}\}+f\_\{\textbackslash\{\}tau\textbackslash\{\}tau\}\textbackslash\{\}nabla\textbackslash\{\}tau\_\textbackslash\{\}psi\textbackslash\{\}otimes\textbackslash\{\}nabla\textbackslash\{\}tau\_\textbackslash\{\}psi+f\_\textbackslash\{\}tau D\_\{\textbackslash\{\}mathbf\{x\}\}\textasciicircum{}2\textbackslash\{\}tau\_\textbackslash\{\}psi,\textbackslash\{\}qquad \textbackslash\{\}Delta f=\textbackslash\{\}operatorname\{tr\}(D\_\{\textbackslash\{\}mathbf\{x\}\}\textasciicircum{}2f)}
\par\emph{为什么：} 只有以正确张量形状保留所有诱导的链式法则项，空间相关的时间坐标才会保持声明的 PDE。脉冲跳变作为界面处理，不对跳变本身求导。
  \item 在脉冲光滑分段内的每个内部点，严格按 E5 求值三类原始物理残差。\(`K_{xi}`\) 必须接收全部声明输入：温度、电场、相态、相态的完整物理梯度和物理拉普拉斯量。严格按 E2，用精确时钟导数与平滑正动力学幅值直接形成 \(`r_{tau}`\)，不得加入对数变换或中心化。使用冻结残差尺度以及非空的初始、边界、本构和脉冲界面项，形成 E5 的物理信息神经网络(Physics-Informed Neural Network, PINN)目标。在脉冲边沿，把同一位置和时刻的左右记录配对，要求温度与相态连续，并在两侧分别求值电学边界条件，使电压可以跟随施加的阶跃。另行积分储热项、导热边界通量和焦耳热来形成能量平衡，只把它作为诊断报告。用 `ResidualBatch` 返回每类残差、权重与标识，以及独立的能量记录。
\noindent\emph{唯一的辅助时钟残差把时钟的精确时间导数直接耦合到平滑的正相变速率幅值。所示归一化相态斜率仅在相态残差与时钟残差同时为零时成立，不构成近似误差界。}
\par\noindent{\small\ttfamily [eq~4, raw LaTeX, not typeset] r\_\textbackslash\{\}tau=\textbackslash\{\}partial\_t\textbackslash\{\}tau\_\textbackslash\{\}psi-\textbackslash\{\}sqrt\{K\_\textbackslash\{\}xi(\textbackslash\{\}hat T,\textbackslash\{\}hat E,\textbackslash\{\}hat\textbackslash\{\}xi,\textbackslash\{\}nabla\textbackslash\{\}hat\textbackslash\{\}xi,\textbackslash\{\}Delta\textbackslash\{\}hat\textbackslash\{\}xi)\textasciicircum{}2+k\_\{\textbackslash\{\}mathrm\{floor\}\}\textasciicircum{}2\},\textbackslash\{\}qquad |\textbackslash\{\}partial\textbackslash\{\}hat\textbackslash\{\}xi/\textbackslash\{\}partial\textbackslash\{\}tau|=|K\_\textbackslash\{\}xi|/\textbackslash\{\}sqrt\{K\_\textbackslash\{\}xi\textasciicircum{}2+k\_\{\textbackslash\{\}mathrm\{floor\}\}\textasciicircum{}2\}<1\textbackslash\{\}quad\textbackslash\{\}text\{when \}r\_\textbackslash\{\}tau=0\textbackslash\{\}text\{ and \}\textbackslash\{\}mathcal\{R\}\_\textbackslash\{\}xi=0}
\noindent\emph{使用完整拉回导数计算原始电学、焦耳热和相态方程，并保留冻结残差尺度以及非空的初值、边界、本构和脉冲界面项；恒等对照令 `tau=t` 时关闭时钟残差项。}
\par\noindent{\small\ttfamily [eq~5, raw LaTeX, not typeset] \textbackslash\{\}mathbf\{R\}\_\{\textbackslash\{\}mathrm\{phys\}\}=\textbackslash\{\}left(\textbackslash\{\}nabla\textbackslash\{\}!\textbackslash\{\}cdot[\textbackslash\{\}sigma(\textbackslash\{\}hat T,\textbackslash\{\}hat\textbackslash\{\}xi)\textbackslash\{\}nabla\textbackslash\{\}hat V],\textbackslash\{\};\textbackslash\{\}rho c\_p\textbackslash\{\}partial\_t\textbackslash\{\}hat T-\textbackslash\{\}nabla\textbackslash\{\}!\textbackslash\{\}cdot[k(\textbackslash\{\}hat T,\textbackslash\{\}hat\textbackslash\{\}xi)\textbackslash\{\}nabla\textbackslash\{\}hat T]-\textbackslash\{\}sigma|\textbackslash\{\}nabla\textbackslash\{\}hat V|\textasciicircum{}2,\textbackslash\{\};\textbackslash\{\}partial\_t\textbackslash\{\}hat\textbackslash\{\}xi-K\_\textbackslash\{\}xi(\textbackslash\{\}hat T,\textbackslash\{\}hat E,\textbackslash\{\}hat\textbackslash\{\}xi,\textbackslash\{\}nabla\textbackslash\{\}hat\textbackslash\{\}xi,\textbackslash\{\}Delta\textbackslash\{\}hat\textbackslash\{\}xi)\textbackslash\{\}right),\textbackslash\{\}quad \textbackslash\{\}mathcal\{L\}=\textbackslash\{\}sum\_q\textbackslash\{\}|\textbackslash\{\}mathcal\{R\}\_q/s\_q\textbackslash\{\}|\_2\textasciicircum{}2+\textbackslash\{\}mathbf\{1\}\_\{\textbackslash\{\}mathrm\{clock\}\}\textbackslash\{\}|r\_\textbackslash\{\}tau/s\_\textbackslash\{\}tau\textbackslash\{\}|\_2\textasciicircum{}2+\textbackslash\{\}mathcal\{L\}\_\{\textbackslash\{\}mathrm\{IC/BC/interface\}\}}
\par\emph{为什么：} 这既保留声明的电-热-相态模型，又使时钟具有动力学特异性。当相态残差和 \(`r_{tau}`\) 都为零时，相态关于时钟时间的斜率幅值小于 `1`；这是条件性质，不是近似保证。
\end{enumerate}
\subsection*{M3\_matched\_training}
\emph{在容量、样本和预算匹配的合同下训练时钟与场。}
\begin{enumerate}[leftmargin=*,start=6]
  \item 建立 `CapacityPlan`，分别统计时钟模块与解模块的可训练参数，只缩减解头的参数分配，使动力学时钟物理信息神经网络(Kinetics-Clock Physics-Informed Neural Network, KC-PINN)总量与声明的参考分配匹配。再建立统一的 `TrainingContract`，其中包含非空的内部、初始、边界、本构和脉冲界面配点标识、冻结残差归一化、从第一次更新起的时钟与场联合更新规则、优化器设置、停止规则，以及单一分布内检查点选择器。 【作者需决定：锁定允许的解头尺寸族与匹配容差、全部优化器设置、更新和停止规则，以及对所有分支统一使用的单一分布内检查点选择标准。】 使用不改变项构成的 第5步（在脉冲光滑分段内的每个内部点） 目标训练，只保存该规则选中的检查点。每个检查点都附带参数数量、配点集标识、尺度标识以及确切的计划和合同版本。
\par\emph{为什么：} 重新分配容量而不是增加容量，可防止参数数量或不同训练合同伪装成时钟效应。
\end{enumerate}
\subsection*{M4\_validation}
\emph{训练匹配基线，运行恒等干预，在分布内锁定全部选择，并在未使用的完整案例上进行分布外评分。}
\begin{enumerate}[leftmargin=*,start=7]
  \item 把每个对照写成冻结的 `BaselineSpec`，并通过 第1步（把现有二维器件模型整理成只读\(\dots \)） 与 第6步（建立 \(`CapacityPl\dots \)） 接口在相同训练---验证案例标识上运行。 【作者需决定：为每个对照提供带版本的实现清单，说明来源、架构或变换、可训练参数、损失项，以及它与共享物理、采样和检查点选择合同的确切映射。】 对照集合必须包括：保留完整拉回和含 \(`K_{xi}`\) 的原相态残差、只移除由 \(`K_{xi}`\) 驱动的时钟量 \(`k_{eff}`\) 与 \(`r_{tau}`\) 的同架构解析单调坐标分支；未修改坐标的多层感知机 PINN；固定解析变换；固定傅里叶与自适应频谱 PINN；因果 PINN 与 NeuSA；自适应和损失条件加权；以及硬约束和 AdamFLIP 对照。另加入实际可训练参数数量与 KC-PINN 相同的加宽未修改坐标分支。所有分支都保留原始物理残差与脉冲界面求值器，并按 第6步（建立 \(`CapacityPl\dots \)） 的统一结构返回检查点。
\par\emph{为什么：} 只有当动力学特异时钟带来的价值超过一般学习坐标、解析时间扭曲、频谱容量、因果性、损失加权和硬约束时，候选方法才成立。
  \item 在同一模型类中使用 `identity` 时钟设置。建立相同的时钟张量和参数对象，将其冻结，并用 `tau\_psi=t` 旁路其输出；同时把 E5 时钟指示量置零，使 \(`r_{tau}`\) 不出现。解头、命名配点集、残差尺度、导数与脉冲界面代码、优化器设置、更新和停止规则以及检查点选择器全部保持不变。同时报告名义参数数量和实际可训练参数数量：休眠时钟参数保证名义匹配，第7步（把每个对照写成冻结的\(\dots \)） 的加宽未修改坐标分支保证实际参数匹配。返回恒等分支检查点和 `ControlTrace`，记录开关状态、旁路确认及共享合同标识。
\par\emph{为什么：} 这个唯一的负对照干预检验后续增益是否需要动力学时钟，而不是来自休眠参数或不同的物理代码。第7步（把每个对照写成冻结的\(\dots \)） 中加宽的未修改坐标分支另行提供实际参数匹配性能。
  \item 在请求的物理时间解码每个检查点。对 KC-PINN，先求值解析 `tau\_psi(x,t)`，再在位置、时钟时间和脉冲值上求值解头；不得反演时钟，也不得沿时钟时间插值。建立 `MetricSpec`，定义场对齐与误差范数、开关事件、相前沿提取与豪斯多夫(Hausdorff)距离、拓扑表示与差异、能量平衡归一化，以及耦合残差的归一化和聚合。 【作者需决定：在最终确定验证评分前，预先声明全部事件阈值、场范数、前沿等值、豪斯多夫与拓扑约定、能量归一化和耦合残差聚合。】 在训练---验证案例上，对 KC-PINN、恒等分支和所有对照应用相同物理时间指标。随后把架构、\(`M_{tau}`\)、时钟基、其他全部设置、阈值、指标代码、停止规则和检查点规则冻结到带版本的 `LockedScoringContract`。
\par\emph{为什么：} 所有估计量都保持为物理时间泛函，而且任何正式分布外结果都不能影响模型或评估选择。
  \item 加载带版本的 `LockedScoringContract` 和从未使用的 `OODCaseManifest`，其中每一行对应一个完整留出几何、材料组合和脉冲协议。逐行加载锁定的 KC-PINN、恒等分支和对照检查点，不更新参数，只构造该案例的物理输入，并独立运行锁定的前向与解码路径。把 第9步（在请求的物理时间解码每个检查\(\dots \)） 的完全相同指标用于物理时间场、开关、前沿、能量平衡和耦合残差。看到任何留出结果后，都不得改变阈值、归一化、架构、时钟基、对照以及停止和检查点规则。按分支和完整案例返回带案例与合同标识的记录，并写明 `evidence\_identity=synthetic\_numerical`；只按锁定规则聚合，不得表述为实验验证。
\par\emph{为什么：} 锁定后的完整实体评分在不泄漏的情况下检验泛化，并把任何正面结论限定为针对所声明 VO2 模型的合成数值证据。
\end{enumerate}
\end{document}
